\section{Transformer解码器的前向计算}
\label{sec:nlp-llm-computation}

第\ref{sec:nlp-llm-data}节把语料变成了定长的词元编号序列。本节固定一个这样的样本$x_{1:T}$，把仅解码器Transformer从输入编号算到下一词元概率的过程逐层写清。要做到可以复算，需要回答四个问题。词元编号只是整数，编号大小没有语义距离，网络需要先把它变成向量；注意力按内容相似度加权求和，本身不携带位置之间的距离，而“周四改到周五”和“周五改到周四”含义相反，网络需要位置信息；一层之内归一化、注意力、前馈与残差怎样连接，张量形状怎样变化；给定一组配置，参数量和每个词元的计算量怎样算出。

第\ref{chap:nlp-classical-models}章的GPT已经采用因果Transformer解码器。今天的大语言模型在这一骨架上做了几处稳定下来的改动：用旋转位置编码代替加到输入上的位置向量，用RMSNorm并放在子层之前，用分组查询注意力减少键值头，用SwiGLU门控前馈层代替ReLU前馈层。下面以LLaMA-3-8B的公开配置为共同对象，把这些部件放进一层完整的计算中，给出层结构图、形状表和参数量手算\scite{nlp-llm-llama3-2024}。注意力、多头、位置编码、RMSNorm与门控前馈层的一般定义和推导可参考\BookChapterRef{chap:transformer}{卷五的Transformer章}；本节只写大语言模型采用的具体形式、配置与核算。第\ref{sec:nlp-llm-long-context}、\ref{sec:nlp-llm-experts}与\ref{sec:nlp-llm-layer-composition}节分别在这条基线上改变注意力读取历史的方式、前馈层调用参数的方式和层间连接的方式。

记号约定如下，并在本章各节沿用。一条序列在第$\ell$层的表示写成矩阵$\bm H^{(\ell)}\in\mathbb R^{T\times d}$，第$t$行$\bm h_t^{(\ell)\top}$是位置$t$的表示，$T$是序列长度，$d$是模型宽度，$L$是层数，$V$是词表大小。$H$表示查询头数，$H_{\mathrm{kv}}$表示键值头数（也就是GQA中的键值组数），$d_h$表示每头维度，通常$Hd_h=d$；它们是标量，与表示矩阵$\bm H$的粗体相区别。前馈中间维度记为$d_{\mathrm{ff}}$，批大小记为$B$。

\subsection{输入嵌入与位置表示}
\label{sec:nlp-llm-embedding}

网络的第一步是给每个编号一个可学习的向量。若直接把编号当作数值输入，编号3与编号4就会被当作“接近”，而这种接近只是词表排序的偶然结果。嵌入表让每个词元拥有自己的一行参数，向量之间的距离由训练决定。第二步是处理顺序。注意力按内容相似度加权求和：在不加掩码的注意力中，把序列位置打乱，各位置的输出只是跟着打乱；因果掩码虽然让各位置看到的前缀长短不同，也没有直接给出两个词元相距多远。模型需要额外的位置信息，才能可靠地区分“周四改到周五”与“周五改到周四”。早期Transformer和GPT把位置向量加到输入嵌入上；当前的大语言模型普遍改用旋转位置编码，它不改动输入嵌入，而是在每一层计算注意力时，把位置写进查询与键的旋转角度。两种做法把位置信息放进了不同的地方，后者让注意力得分只依赖两个位置的相对距离。

\subsubsection{词元嵌入与输出层权重共享}
\label{sec:nlp-llm-embedding-tying}

设词表大小为$V$，嵌入矩阵为$\bm E\in\mathbb R^{V\times d}$。位置$t$的输入表示是该词元对应的那一行：
\[
  \bm h_t^{(0)}=\bm E_{x_t,:}^{\top}\in\mathbb R^{d},
  \qquad
  \bm H^{(0)}=\bigl[\bm E_{x_1,:};\ \ldots;\ \bm E_{x_T,:}\bigr]\in\mathbb R^{T\times d}.
\]
取行等价于用独热向量乘以$\bm E$，但实现中只做查表，不计入矩阵乘法的计算量。同一词元在任何位置、任何样本中读取同一行参数，所以一处样本对这一行的更新会影响所有含该词元的样本；同一词元在不同句子中的不同含义，要靠后续各层的上下文计算区分。

在输出端，网络把最后一层的表示$\bm z_t\in\mathbb R^d$映射为词表上的$V$个得分（logits），即$\bm W_{\mathrm{out}}\bm z_t$，其中$\bm W_{\mathrm{out}}\in\mathbb R^{V\times d}$。它与嵌入矩阵形状相同，每一行也对应一个词元。\term{参数绑定}（weight tying）令$\bm W_{\mathrm{out}}=\bm E$：词元$v$的输出得分就是当前表示与它的输入嵌入的内积，输入端和输出端共享同一组词元向量，省去$Vd$个参数\scite{nlp-llm-press2017tying}。直觉上，在输入端出现于相似上下文的词元，在输出端也应在相似上下文中被预测，共享参数让两端的训练信号互相补充，稀有词元尤其受益。

是否共享取决于嵌入在总参数中占多大比例。GPT-2的小型配置$d=768$、$V=50257$，嵌入矩阵约有3860万个参数，约占全部参数的三成，它采用了参数绑定\scite{radford2019gpt2}。LLaMA-3-8B的$V=128256$、$d=4096$，一个嵌入矩阵约5.25亿个参数，只占约80亿总参数的6.5\%，它没有共享输出层。参数规模增大后，共享节省的比例下降，而独立的输出矩阵让输入表示与输出分类各自优化，不受“同一组向量兼顾两种用途”的约束。两种选择在LLaMA-3-8B上相差约5.25亿个参数，第\ref{sec:nlp-llm-decoder-cost}小节会把这项计入总量。

\subsubsection{旋转位置编码的多频率形式}
\label{sec:nlp-llm-rope}

\term{旋转位置编码}（rotary position embedding, RoPE）在二维情形下的思想是：把位置$t$的查询和位置$j$的键分别旋转角度$t\omega$与$j\omega$，二者点积中的角度只剩差值$(j-t)\omega$\scite{su2021roformer}。其推导只用到旋转矩阵的两条性质：转置等于反向旋转，两次旋转的复合等于角度相加。下面写出高维向量上使用的多频率形式。

设某一头的查询或键为$\bm u\in\mathbb R^{d_h}$，$d_h$为偶数。把坐标两两配对成$d_h/2$个二维对，第$i$对（$i=0,1,\ldots,d_h/2-1$）使用角频率
\[
  \theta_i=b^{-2i/d_h},
\]
其中$b$是底数，原始RoPE取$b=10000$。位置$t$上的旋转矩阵是分块对角矩阵$\bm R_t=\operatorname{diag}\bigl(\bm R(t\theta_0),\ldots,\bm R(t\theta_{d_h/2-1})\bigr)$，每个块是二维旋转$\bm R(\phi)=\bigl(\begin{smallmatrix}\cos\phi&-\sin\phi\\ \sin\phi&\cos\phi\end{smallmatrix}\bigr)$。由各块的复合关系$\bm R(t\theta_i)^{\top}\bm R(j\theta_i)=\bm R((j-t)\theta_i)$，旋转后的点积为
\begin{equation}
  (\bm R_t\bm q_t)^{\top}(\bm R_j\bm k_j)
  =\bm q_t^{\top}\bm R_{j-t}\bm k_j
  =\sum_{i=0}^{d_h/2-1}\Bigl[c_i\cos\bigl((j-t)\theta_i\bigr)+s_i\sin\bigl((j-t)\theta_i\bigr)\Bigr],
  \label{eq:nlp-llm-rope}
\end{equation}
其中$c_i$与$s_i$是第$i$对坐标上查询与键的两种乘积组合（分别为对应坐标乘积之和与交叉乘积之差），只由内容决定。位置只以差值$j-t$出现：内容不变时，把查询和键同时平移，得分不变。RoPE只作用于查询和键，值向量不旋转；实现中也有把前后两半坐标而不是相邻坐标配对的写法，两者相差一个固定的坐标排列，性质相同。

多频率的意义在于让不同坐标对负责不同尺度的距离。第$i$对的旋转波长是$\lambda_i=2\pi/\theta_i$，即$2\pi b^{2i/d_h}$个位置。取$d_h=128$：$i=0$时$\lambda_0\approx6.3$，相邻几个位置的角度差就很明显，这一对对短距离的变化敏感；$i$增大时波长按几何级数变长，最后一对的角度在很长的距离上才转完一圈，用来区分远距离。$b=10000$时最长波长约为5.4万个位置；LLaMA-3把底数提高到$b=500000$，最长波长约为256万个位置，$i=32$的中间一对也已达到约4400个位置\scite{nlp-llm-llama3-2024}。

这一结构也说明了长度外推的困难。若训练序列长度为$T_{\mathrm{train}}$，波长大于$T_{\mathrm{train}}$的低频对在训练中只见过不足一圈的角度范围。推断时若距离超过$T_{\mathrm{train}}$，这些坐标对会出现训练中从未见过的角度组合，注意力得分的分布随之偏移；高频对则早已在训练中转过许多圈，不存在这个问题。式\eqref{eq:nlp-llm-rope}在任意位置都可以求值，但这不等于模型在任意长度上都能正确使用它。第\ref{sec:nlp-llm-context-extension}小节的位置插值、底数调整等方法，都是在改造这些低频对的角度范围。

把RoPE放回位置方案的谱系中比较，可以看清它解决了什么。GPT-2使用可学习的绝对位置表，每个位置一行向量加到输入嵌入上，表的行数（1024）就是能处理的最大长度，超出部分没有参数可用。固定的正弦位置向量不需要学习，但同样是加到输入上，位置与内容在之后的每个投影中混在一起，注意力得分中出现内容与位置的交叉项，相对距离只能间接得到。ALiBi不使用位置向量，而是在注意力得分上直接减去与距离成正比的偏置，结构简单、外推较平稳，但把“远处的词元权重更低”写死在得分里\scite{press2021alibi}。RoPE在每一层的查询与键上施加旋转，不需要位置参数表，得分中的位置只以相对距离出现，并且仍与内容相乘而不是简单相加。它的代价是每层多一次逐坐标的旋转运算，以及上面所说的低频对外推问题。

\subsection{一层解码器的完整计算}
\label{sec:nlp-llm-decoder-layer}

有了输入表示$\bm H^{(0)}$，网络把同一结构的解码器层重复$L$次。图~\ref{fig:nlp-llm-decoder-layer}画出了一层的全部运算。读图时可以沿主链路自上而下：输入先经RMSNorm，再分三路投影为查询、键和值，查询与键施加RoPE，各头做因果注意力，拼接后经输出投影，与输入残差相加；随后经第二个RMSNorm，进入SwiGLU前馈层，再与前馈层的输入残差相加。右侧两条旁路说明每个子层只计算对表示的修改量，左侧灰字给出各处的张量形状。下面按图中的顺序逐个写出运算，最后汇总为公式链与形状表。

\begin{figure}[htbp]
  \centering
  % 第九卷 3.2.2 节：前置 RMSNorm 的仅解码器 Transformer 层（GQA 注意力 + SwiGLU 前馈）。
% 主链路自上而下；右侧灰线为两条残差旁路；左侧灰字为单条序列（省略批维）上的张量形状。
% 蓝色：注意力子层；黄色：前馈子层；白底灰框：归一化与逐元素运算。
\begin{tikzpicture}[
  x=1mm,y=1mm,
  every node/.style={font=\figurefont\footnotesize,text=FigureInk},
  op/.style={rectangle,rounded corners=1.8pt,draw=FigureStroke,line width=0.6pt,fill=white,
    minimum height=6.5mm,inner sep=1.2mm,align=center},
  attnop/.style={op,draw=FigureBlue,fill=FigureBlueFill},
  ffnop/.style={op,draw=FigureYellow,fill=FigureYellowFill},
  dot/.style={circle,draw=FigureStroke,line width=0.6pt,fill=white,inner sep=0pt,minimum size=4.6mm},
  flow/.style={-{Stealth[length=1.8mm,width=1.3mm]},draw=FigureStroke,line width=0.7pt},
  skip/.style={-{Stealth[length=1.8mm,width=1.3mm]},draw=FigureMuted,line width=0.7pt,rounded corners=1.8pt},
  dim/.style={font=\figurefont\scriptsize,text=FigureMuted,inner sep=0.8pt},
  panel/.style={rounded corners=2.5pt,line width=0.5pt},
  ptitle/.style={font=\figurefont\scriptsize\bfseries,anchor=north west}
]
  % 子层分组背景。
  \fill[panel,FigureBlueFill!45] (-6,-10) rectangle (76,-93);
  \fill[panel,FigureYellowFill!55] (-6,-97) rectangle (76,-166);
  \node[ptitle,text=FigureBlue] at (-5,-11) {注意力子层};
  \node[ptitle,text=FigureYellow!80!black] at (-5,-98) {前馈子层};

  % 输入与第一个归一化。
  \node[op,draw=none,fill=none] (in) at (35,0) {层输入$\bm H^{(\ell)}$};
  \node[op,minimum width=24mm] (n1) at (35,-18) {RMSNorm};
  \draw[flow] (in) -- (n1);
  \node[dim,anchor=east] at (33,-8) {$T\times d$};

  % 三个投影。
  \node[attnop,minimum width=13mm] (wq) at (15,-32) {$\bm W_Q$};
  \node[attnop,minimum width=13mm] (wk) at (35,-32) {$\bm W_K$};
  \node[attnop,minimum width=13mm] (wv) at (55,-32) {$\bm W_V$};
  \coordinate (split) at (35,-25);
  \draw[draw=FigureStroke,line width=0.7pt] (n1.south) -- (split);
  \draw[flow] (split) -| (wq.north);
  \draw[flow] (split) -- (wk.north);
  \draw[flow] (split) -| (wv.north);

  % RoPE 只作用于查询和键。
  \node[op,minimum width=13mm] (rq) at (15,-46) {RoPE};
  \node[op,minimum width=13mm] (rk) at (35,-46) {RoPE};
  \draw[flow] (wq) -- (rq);
  \draw[flow] (wk) -- (rk);
  \node[dim,anchor=east] at (13.5,-39) {$T\times H\times d_h$};
  \node[dim,anchor=west] at (36.5,-39) {$T\times H_{kv}\times d_h$};
  \node[dim,anchor=west] at (56.5,-46) {$T\times H_{kv}\times d_h$};

  % 因果注意力与输出投影。
  \node[attnop,minimum width=62mm] (att) at (35,-62)
    {每头因果注意力\quad$\operatorname{softmax}\bigl(\bm Q_a\bm K_{g(a)}^{\top}/\sqrt{d_h}+\bm M\bigr)\bm V_{g(a)}$};
  \draw[flow] (rq) -- (rq |- att.north);
  \draw[flow] (rk) -- (rk |- att.north);
  \draw[flow] (wv) -- (wv |- att.north);
  \node[dim,anchor=north east] at (att.south east) {得分：每头$T\times T$};
  \node[attnop,minimum width=40mm] (wo) at (35,-77) {拼接各头，乘$\bm W_O$};
  \draw[flow] (att) -- (wo);
  \node[dim,anchor=east] at (33,-70.5) {$T\times Hd_h$};
  \node[dot] (a1) at (35,-88) {$+$};
  \draw[flow] (wo) -- (a1);

  % 第一条残差旁路。
  \coordinate (r1) at (35,-7);
  \draw[skip] (r1) -- (72,-7) -- (72,-88) -- (a1.east);

  % 前馈子层。
  \node[op,minimum width=24mm] (n2) at (35,-105) {RMSNorm};
  \draw[flow] (a1) -- (n2);
  \node[dim,anchor=east] at (33,-94.5) {$T\times d$};
  \node[ffnop,minimum width=13mm] (w1) at (22,-119) {$\bm W_1$};
  \node[ffnop,minimum width=13mm] (w3) at (48,-119) {$\bm W_3$};
  \coordinate (split2) at (35,-112);
  \draw[draw=FigureStroke,line width=0.7pt] (n2.south) -- (split2);
  \draw[flow] (split2) -| (w1.north);
  \draw[flow] (split2) -| (w3.north);
  \node[op,minimum width=13mm] (silu) at (22,-132) {SiLU};
  \draw[flow] (w1) -- (silu);
  \node[dim,anchor=east] at (20.5,-125.5) {$T\times d_{\mathrm{ff}}$};
  \node[dot] (m) at (35,-142) {$\odot$};
  \draw[flow] (silu) |- (m.west);
  \draw[flow] (w3) |- (m.east);
  \node[dim,anchor=west] at (49.5,-130) {$T\times d_{\mathrm{ff}}$};
  \node[ffnop,minimum width=13mm] (w2) at (35,-153) {$\bm W_2$};
  \draw[flow] (m) -- (w2);
  \node[dot] (a2) at (35,-162) {$+$};
  \draw[flow] (w2) -- (a2);

  % 第二条残差旁路。
  \coordinate (r2) at (35,-94.5);
  \draw[skip] (r2) -- (72,-94.5) -- (72,-162) -- (a2.east);

  % 输出。
  \node[op,draw=none,fill=none] (out) at (35,-176) {层输出$\bm H^{(\ell+1)}$};
  \draw[flow] (a2) -- (out);
  \node[dim,anchor=east] at (33,-169.5) {$T\times d$};
  \node[dim,anchor=west,text=FigureMuted] at (73,-50) {残差};
  \node[dim,anchor=west,text=FigureMuted] at (73,-128) {残差};
\end{tikzpicture}

  \caption{一层前置归一化的解码器层。蓝色区域为注意力子层：查询、键、值投影后只对查询和键施加RoPE，第$a$个查询头读取第$g(a)$组键值，$\bm M$为因果掩码；各头输出拼接后由$\bm W_O$映射回$d$维。黄色区域为SwiGLU前馈子层，$\bm W_1$分支经SiLU后与$\bm W_3$分支逐元素相乘，再由$\bm W_2$映射回$d$维。右侧灰线是残差旁路，左侧灰字是省略批维后的张量形状。}
  \label{fig:nlp-llm-decoder-layer}
\end{figure}

\subsubsection{RMSNorm与前置归一化}
\label{sec:nlp-llm-rmsnorm}

各层的输出逐层累加到残差流上，数值尺度会随深度漂移。进入每个子层之前先归一化，可以让子层看到尺度稳定的输入。\term{均方根归一化}（root mean square normalization, RMSNorm）对每个位置的表示$\bm x\in\mathbb R^d$单独计算
\begin{equation}
  \operatorname{RMSNorm}(\bm x)=\frac{\bm x}{\sqrt{\frac1d\sum_{i=1}^{d}x_i^2+\varepsilon}}\odot\bm g,
  \label{eq:nlp-llm-rmsnorm}
\end{equation}
其中$\bm g\in\mathbb R^d$是可学习的逐坐标尺度，$\varepsilon$是防止除零的小常数，LLaMA-3取$10^{-5}$\scite{zhang2019rmsnorm}。与层归一化相比，它不减去均值，也没有偏置，少一次求均值和一次减法。例如$\bm x=(1,-1,3,1)$，均方为$(1+1+9+1)/4=3$，取$\bm g=\bm 1$并忽略$\varepsilon$，结果为$(0.58,-0.58,1.73,0.58)$；层归一化会先减去均值1，得到$(0,-1.41,1.41,0)$。两者都消除了整体缩放，只有层归一化消除了整体平移。

归一化放在哪里同样重要。大语言模型普遍采用\term{前置层归一化}（pre-layer normalization, Pre-LN），每个子层写成$\bm x+F(\operatorname{Norm}(\bm x))$：归一化只作用于进入子层$F$的分支，残差旁路上的$\bm x$原样传下去。最后一层之后再接一个最终归一化，然后送入输出层。前置与后置两种位置在梯度传播和训练稳定性上的比较，在第\ref{sec:nlp-llm-residual-normalization}小节展开。

\subsubsection{因果多头注意力}
\label{sec:nlp-llm-causal-mha}

注意力子层的输入是归一化后的$\bm X=\operatorname{RMSNorm}(\bm H^{(\ell)})\in\mathbb R^{T\times d}$。三个投影矩阵把它映射为查询、键和值：
\[
  \bm Q=\bm X\bm W_Q,\quad \bm K=\bm X\bm W_K,\quad \bm V=\bm X\bm W_V,
\]
其中$\bm W_Q\in\mathbb R^{d\times Hd_h}$，$\bm W_K,\bm W_V\in\mathbb R^{d\times H_{\mathrm{kv}}d_h}$。
把$\bm Q$的列每$d_h$个分为一组，得到$H$个查询头$\bm Q_a\in\mathbb R^{T\times d_h}$；$\bm K$与$\bm V$同样分为$H_{\mathrm{kv}}$个键值头。查询头多于键值头时，\term{分组查询注意力}（grouped-query attention, GQA）让每$H/H_{\mathrm{kv}}$个查询头共用一组键值，第$a$个查询头使用第$g(a)=\lceil aH_{\mathrm{kv}}/H\rceil$组\scite{ainslie2023gqa}。$H_{\mathrm{kv}}=H$时就是标准多头注意力。LLaMA-3-8B的$H=32$、$H_{\mathrm{kv}}=8$，每4个查询头共用一组键值，这把推断时需要缓存的键值减为四分之一；它的动机、由多头检查点转换的方法以及其他压缩方式在第\ref{sec:nlp-llm-kv-compression}小节讨论。投影之后，对每个查询头和键头的每个位置施加式\eqref{eq:nlp-llm-rope}中的旋转。

对第$a$个头，省略头编号，记$\bm q_t$、$\bm k_j$、$\bm v_j$为旋转后的查询、旋转后的键与值在相应位置的行（写成列向量）。\term{因果注意力}（causal attention）只允许位置$t$读取不晚于它的位置：
\begin{equation}
  \alpha_{t,j}=\frac{\exp\bigl(\bm q_t^{\top}\bm k_j/\sqrt{d_h}\bigr)}
  {\sum_{i=1}^{t}\exp\bigl(\bm q_t^{\top}\bm k_i/\sqrt{d_h}\bigr)}\quad(j\leq t),
  \qquad
  \bm o_t=\sum_{j=1}^{t}\alpha_{t,j}\bm v_j.
  \label{eq:nlp-llm-causal-attention}
\end{equation}
写成矩阵形式，令掩码$\bm M\in\mathbb R^{T\times T}$在$j\leq t$处为0、在$j>t$处为$-\infty$，则$\bm O_a=\operatorname{softmax}\bigl(\bm Q_a\bm K_{g(a)}^{\top}/\sqrt{d_h}+\bm M\bigr)\bm V_{g(a)}$，softmax逐行计算。$-\infty$经过指数变为0，未来位置的权重恰好为零，这保证了训练时一次前向同时计算所有位置，也不会让任何位置看到它要预测的词元。除以$\sqrt{d_h}$控制点积随维度增长的尺度。

各头独立计算后，\term{多头注意力}（multi-head attention, MHA）把$H$个输出沿特征维拼接，再由输出投影混合：
\begin{equation}
  \operatorname{Attn}(\bm X)=\bigl[\bm O_1,\ \bm O_2,\ \ldots,\ \bm O_H\bigr]\bm W_O,
  \qquad \bm W_O\in\mathbb R^{Hd_h\times d}.
  \label{eq:nlp-llm-mha}
\end{equation}
拼接保留了各头的结果，输出投影再把它们线性组合回$d$维。哪个头形成什么样的读取模式由训练决定，并没有预先指定。

\Needspace{10\baselineskip}
\begin{example}[两头因果注意力的手算]
  \label{ex:nlp-llm-two-head-attention}
  取$T=3$、$H=H_{\mathrm{kv}}=2$、$d_h=2$、$d=4$，数值为教学构造，假设投影与RoPE已经完成。第一个头的查询、键与值相同，三行依次为$(1,0)$、$(0,1)$、$(1,1)$。缩放后的得分矩阵$\bm Q_1\bm K_1^{\top}/\sqrt2$的下三角部分为
  \[
    \begin{pmatrix}
      0.71&-\infty&-\infty\\
      0&0.71&-\infty\\
      0.71&0.71&1.41
    \end{pmatrix},
  \]
  其中第2行第3列的原始得分是0.71，被掩码替换为$-\infty$。逐行softmax得到权重$(1)$、$(0.33,0.67)$、$(0.25,0.25,0.50)$，乘以值得到$\bm O_1$的三行$(1,0)$、$(0.33,0.67)$、$(0.75,0.75)$。

  第二个头的三个查询都是$(0,1)$，键为$(0,2)$、$(0,0)$、$(0,0)$，值为$(2,0)$、$(0,2)$、$(0,0)$。每个位置对第1个位置的得分为$2/\sqrt2\approx1.41$，对其余位置为0，权重依次为$(1)$、$(0.80,0.20)$、$(0.67,0.16,0.16)$，输出为$(2,0)$、$(1.61,0.39)$、$(1.35,0.33)$。

  两头输出沿特征维拼接，得到$3\times4$矩阵
  \[
    \bigl[\bm O_1,\ \bm O_2\bigr]=
    \begin{pmatrix}
      1&0&2&0\\
      0.33&0.67&1.61&0.39\\
      0.75&0.75&1.35&0.33
    \end{pmatrix}.
  \]
  若$\bm W_O=\bm I_4$，这就是注意力子层的输出，再与该层输入相加；一般的$\bm W_O$会把两个头的四个坐标混合。两个头在同一输入上形成了不同的读取：第一个头按内容匹配，位置3把一半权重给自己；第二个头的键只在位置1上非零，各位置都主要读取开头。位置1在两个头中都只能读自己，这是掩码而不是得分决定的。
\end{example}

\subsubsection{SwiGLU门控前馈层}
\label{sec:nlp-llm-swiglu}

注意力把其他位置的信息汇集到当前位置，但它对值向量只做加权平均，是线性的。前馈层在每个位置上独立做非线性变换。原始Transformer的前馈层是$\max(0,\bm x\bm W_1)\bm W_2$，中间维度为$4d$。大语言模型多采用门控形式的\term{逐位置前馈网络}（position-wise feed-forward network, FFN）。对归一化后的输入$\bm Y=\operatorname{RMSNorm}(\bm U)\in\mathbb R^{T\times d}$，
\begin{equation}
  \operatorname{FFN}(\bm Y)=\Bigl[\operatorname{SiLU}(\bm Y\bm W_1)\odot(\bm Y\bm W_3)\Bigr]\bm W_2,
  \label{eq:nlp-llm-ffn}
\end{equation}
其中$\bm W_1,\bm W_3\in\mathbb R^{d\times d_{\mathrm{ff}}}$，$\bm W_2\in\mathbb R^{d_{\mathrm{ff}}\times d}$；$\operatorname{SiLU}(z)=z\sigma(z)$逐元素作用，$\sigma$是sigmoid函数，$\odot$是逐元素乘法。这是\term{门控线性单元}（gated linear unit, GLU）的SwiGLU变体\scite{shazeer2020glu}。两路投影的作用不同：$\bm Y\bm W_3$提供候选特征，$\operatorname{SiLU}(\bm Y\bm W_1)$按输入内容逐坐标放大、缩小或翻转这些特征，乘法让一路投影调节另一路。前馈层的计算只涉及同一行，不同位置之间没有交互，这与注意力形成分工。Shazeer在相同参数量与计算量的条件下比较了多种前馈层，门控变体在预训练困惑度和若干下游任务上优于ReLU与GELU前馈层；该文没有给出机制上的解释，这一选择主要依据经验比较\scite{shazeer2020glu}。代价是多一个投影矩阵和一次逐元素乘法，所以中间维度要相应缩小，才能保持参数量不变。

门控前馈层有三个矩阵，参数为$3d\,d_{\mathrm{ff}}$；两矩阵ReLU前馈层在$4d$中间维度下有$8d^2$个参数。令二者相等得$d_{\mathrm{ff}}=8d/3$，这是常见的取值起点。实现中再乘一个倍率并向上取整到硬件友好的倍数：LLaMA-3-8B从$8\times4096/3\approx10923$出发乘以1.3，再向上取整到1024的倍数，得到$d_{\mathrm{ff}}=14336=3.5d$。

\subsubsection{残差连接与张量形状}
\label{sec:nlp-llm-residual-shapes}

把前三个部件按图~\ref{fig:nlp-llm-decoder-layer}连起来，第$\ell$层的完整计算是
\begin{equation}
  \begin{aligned}
    \bm U^{(\ell)}&=\bm H^{(\ell)}+\operatorname{Attn}\bigl(\operatorname{RMSNorm}(\bm H^{(\ell)})\bigr),\\
    \bm H^{(\ell+1)}&=\bm U^{(\ell)}+\operatorname{FFN}\bigl(\operatorname{RMSNorm}(\bm U^{(\ell)})\bigr),
  \end{aligned}
  \qquad \ell=0,1,\ldots,L-1.
  \label{eq:nlp-llm-decoder-layer}
\end{equation}
两个RMSNorm各有自己的尺度参数。$L$层之后，最终表示为$\bm Z=\operatorname{RMSNorm}(\bm H^{(L)})$，第$t$行$\bm z_t$送入输出层。\term{残差连接}（residual connection）让每个子层只输出对表示的修改量：若某层暂时学不到有用的变换，它可以输出接近零的修改，而不破坏前面各层已经形成的表示。残差流$\bm H^{(\ell)}$的形状在整个网络中保持$T\times d$不变，各子层从它读出、再写回，这也使得后几节改动注意力或前馈层时，层间接口可以保持不变。

表~\ref{tab:nlp-llm-decoder-shapes}列出一次前向中主要张量的形状及其在LLaMA-3-8B上的取值。有两处值得留意。其一是每头的$T\times T$得分矩阵：$T=8192$、$H=32$时，一层的得分元素有$32\times8192^2\approx2.1\times10^{9}$个，若全部写入显存，会远超权重本身所占的空间。实际实现按块计算softmax，不保存完整的得分矩阵，这一访存优化在第\ref{chap:nlp-model-inference}章讨论。其二是前馈中间张量，宽度$d_{\mathrm{ff}}$是$d$的3.5倍，是逐位置计算中最宽的部分。

\begin{table}[htbp]
  \centering
  \begin{tabular}{@{}lll@{}}
    \toprule
    张量 & 形状 & LLaMA-3-8B取值\\
    \midrule
    输入编号 & $B\times T$ & ---\\
    残差流$\bm H^{(\ell)}$、$\bm U^{(\ell)}$ & $B\times T\times d$ & $d=4096$\\
    查询$\bm Q$ & $B\times T\times H\times d_h$ & $32\times128$\\
    键$\bm K$、值$\bm V$ & $B\times T\times H_{\mathrm{kv}}\times d_h$ & $8\times128$\\
    得分与权重 & $B\times H\times T\times T$ & 每头$T\times T$\\
    拼接后的注意力输出 & $B\times T\times Hd_h$ & $4096$\\
    前馈中间张量 & $B\times T\times d_{\mathrm{ff}}$ & $14336$\\
    输出得分 & $B\times T\times V$ & $128256$\\
    \bottomrule
  \end{tabular}
  \caption{解码器一次前向中主要张量的形状。$B$为批大小，$T$为序列长度；GQA下键和值的头数$H_{\mathrm{kv}}$小于查询头数$H$。注意力得分矩阵随$T$平方增长，其余张量随$T$线性增长。}
  \label{tab:nlp-llm-decoder-shapes}
\end{table}

这些部件并非一开始就是这样组合的。表~\ref{tab:nlp-llm-decoder-choices}把第\ref{chap:nlp-classical-models}章介绍的GPT系列中的GPT-2与LLaMA-3-8B逐项对照。两者都是前置归一化的因果解码器，层的骨架相同；差别集中在位置信息进入的方式、归一化的形式、前馈层的形式、键值头的数量和输出层是否共享。每一项替换都只改变一个部件，接口仍是$T\times d$的残差流，这正是它们可以分别比较、逐步采用的原因。

\begin{table*}[htbp]
  \begin{tabularx}{\linewidth}{@{}p{24mm}p{40mm}p{40mm}X@{}}
    \toprule
    部件 & GPT-2 & LLaMA-3-8B & 改变的作用\\
    \midrule
    位置信息 & 可学习绝对位置表，加到输入 & 每层对查询与键施加RoPE & 得分只依赖相对距离；无位置参数表\\
    归一化 & 前置LayerNorm，含偏置 & 前置RMSNorm，无偏置 & 少一次求均值；不消除平移\\
    前馈层 & GELU，中间维度$4d$ & SwiGLU，中间维度$3.5d$ & 门控调节特征；多一个矩阵\\
    注意力头 & 标准多头，$H_{\mathrm{kv}}=H$ & GQA，$H_{\mathrm{kv}}=H/4$ & 键值投影与缓存减为四分之一\\
    输出层 & 与嵌入共享 & 独立矩阵 & 嵌入占比小时放开两端约束\\
    线性层偏置 & 有 & 无 & 参数与实现略简\\
    \bottomrule
  \end{tabularx}
  \caption{GPT-2与LLaMA-3-8B解码器层的部件对照。两者的层骨架相同，都是带残差的前置归一化子层；表中只列出被替换的部件及替换后改变的对象，各项收益依赖相应的经验比较。}
  \label{tab:nlp-llm-decoder-choices}
\end{table*}

\subsection{参数量、计算量与典型配置}
\label{sec:nlp-llm-decoder-cost}

模型规模通常用参数量描述，训练和推断的成本则取决于每个词元的计算量。两者都可以从配置直接推出。先数一层的参数。注意力子层有四个投影矩阵，$\bm W_Q$与$\bm W_O$各有$d\cdot Hd_h$个参数，$\bm W_K$与$\bm W_V$各有$d\cdot H_{\mathrm{kv}}d_h$个；前馈子层有三个矩阵；两个RMSNorm各有$d$个尺度参数。在$Hd_h=d$时，
\begin{equation}
  P_{\mathrm{attn}}=2d^2\Bigl(1+\frac{H_{\mathrm{kv}}}{H}\Bigr),\qquad
  P_{\mathrm{ffn}}=3d\,d_{\mathrm{ff}},\qquad
  P_{\mathrm{layer}}=P_{\mathrm{attn}}+P_{\mathrm{ffn}}+2d.
  \label{eq:nlp-llm-layer-params}
\end{equation}
标准多头注意力$H_{\mathrm{kv}}=H$时$P_{\mathrm{attn}}=4d^2$，GQA把键值投影缩小，注意力参数介于$2d^2$与$4d^2$之间。全模型的参数量为$N=L\,P_{\mathrm{layer}}+Vd+d$，若输出层不共享，再加$Vd$。

表~\ref{tab:nlp-llm-llama3-config}列出LLaMA-3-8B的配置。代入式\eqref{eq:nlp-llm-layer-params}逐项计算：
\begin{itemize}
  \item 注意力：$2\times4096^2\times(1+8/32)=41{,}943{,}040$，约0.42亿；
  \item 前馈：$3\times4096\times14336=176{,}160{,}768$，约1.76亿；
  \item 每层合计（含两个RMSNorm的8192个参数）：$218{,}112{,}000$，约2.18亿，其中前馈约占81\%；
  \item 32层：$6{,}979{,}584{,}000$，约69.8亿；
  \item 嵌入与不共享的输出层：$2\times128256\times4096=1{,}050{,}673{,}152$，约10.5亿；
  \item 再加最终RMSNorm的4096个参数，总计$8{,}030{,}261{,}248$，约80.3亿。
\end{itemize}
这与“8B”的名称一致。对标准多头注意力与$d_{\mathrm{ff}}=4d$的两矩阵前馈层，每层参数为$4d^2+8d^2=12d^2$，常用$N\approx12Ld^2$粗估非嵌入参数；GQA与SwiGLU改变了其中的系数，LLaMA-3-8B的每层约为$13d^2$。可以看到，前馈层而不是注意力占了参数的大头；嵌入和输出层合计约占13\%，在这个规模上已不可忽略，这也解释了为什么词表大小是需要权衡的设计。

\begin{table}[htbp]
  \centering
  \begin{tabular}{@{}llr@{}}
    \toprule
    配置项 & 记号 & 取值\\
    \midrule
    模型宽度 & $d$ & 4096\\
    层数 & $L$ & 32\\
    查询头数 & $H$ & 32\\
    键值头数 & $H_{\mathrm{kv}}$ & 8\\
    每头维度 & $d_h$ & 128\\
    前馈中间维度 & $d_{\mathrm{ff}}$ & 14336\\
    词表大小 & $V$ & 128256\\
    RoPE底数 & $b$ & 500000\\
    RMSNorm常数 & $\varepsilon$ & $10^{-5}$\\
    输出层与嵌入共享 & --- & 否\\
    \bottomrule
  \end{tabular}
  \caption{LLaMA-3-8B的结构配置，取自该模型的公开报告与发布配置。主阶段预训练的序列长度为8192，后续继续预训练把上下文扩展到128K。}
  \label{tab:nlp-llm-llama3-config}
\end{table}

再看计算量。一次矩阵乘法$\bm y=\bm W\bm x$中，$\bm W$的每个元素对应一次乘法和一次加法，即2次\term{浮点运算}（floating-point operations, FLOPs）。一个词元前向经过所有投影矩阵、前馈矩阵和输出层，每个参数恰好参与一次乘加，因此每词元的前向计算量约为
\begin{equation}
  C_{\mathrm{fwd}}\approx2N_{\mathrm{mat}}+2LTd,
  \label{eq:nlp-llm-forward-flops}
\end{equation}
其中$N_{\mathrm{mat}}$是参与矩阵乘法的参数数，不含嵌入查表与归一化尺度\scite{kaplan2020scaling}。第二项来自注意力中不含参数的两步乘法：位置$t$的一个查询头与$t$个键做点积，需要$t\,d_h$次乘加，$H$个头共$t\,d$次；按权重对$t$个值求和又是$t\,d$次。一层合计$2t\,d$次乘加即$4t\,d$次浮点运算。在因果序列中$t$从1取到$T$，平均约为$T/2$，于是平均每词元每层约$2Td$次，$L$层合计$2LTd$。注意GQA减少的是键值投影和缓存，查询头仍是$H$个，这一项并不随之减小。

以LLaMA-3-8B为例，参与矩阵乘法的参数是32层的注意力与前馈矩阵再加独立的输出层（嵌入查表不计）：
\[
  N_{\mathrm{mat}}=32\times(41{,}943{,}040+176{,}160{,}768)+525{,}336{,}576\approx7.50\times10^{9},
\]
参数项约为$1.50\times10^{10}$次浮点运算，即每词元约150亿次。$T=8192$时，注意力项为$2\times32\times8192\times4096\approx2.1\times10^{9}$，约为参数项的14\%；$T=131072$时增长到约$3.4\times10^{10}$，是参数项的2.3倍。短上下文下每词元计算量近似为$2N$，长上下文下注意力项变为主导；推断时还要为每个历史词元保存键和值。这两项资源代价在第\ref{sec:nlp-llm-context-bottlenecks}小节一并核算。

训练还需要反向传播。对每个矩阵乘法$\bm Y=\bm X\bm W$，反向既要计算对输入的梯度$\partial\mathcal L/\partial\bm X=(\partial\mathcal L/\partial\bm Y)\bm W^{\top}$以便继续向前传，又要计算对权重的梯度$\partial\mathcal L/\partial\bm W=\bm X^{\top}(\partial\mathcal L/\partial\bm Y)$，两次乘法的规模都与前向相同，所以反向约为前向的两倍，即每词元约$4N$次。前向与反向相加，每个训练词元的计算量约为$6N$，训练$D$个词元的总计算量为
\begin{equation}
  C\approx6ND.
  \label{eq:nlp-llm-train-flops}
\end{equation}
式\eqref{eq:nlp-llm-train-flops}略去了注意力项，常规长度下误差在一成多以内，长上下文阶段需要另计；稀疏专家模型中$N$应取每词元的激活参数量。这一估计怎样用于在固定预算下选择模型与数据规模，在第\ref{sec:nlp-llm-compute-allocation}小节展开。

\subsection{词元概率与序列概率}
\label{sec:nlp-llm-probabilities}

经过$L$层与最终归一化，位置$t$得到$\bm z_t\in\mathbb R^d$。输出层给词表中每个候选一个得分，softmax把得分归一化为下一词元的分布：
\begin{equation}
  p_\theta(x_{t+1}=v\mid x_{\leq t})=
  \frac{\exp\bigl((\bm W_{\mathrm{out}}\bm z_t+\bm c)_v\bigr)}
  {\sum_{u=1}^{V}\exp\bigl((\bm W_{\mathrm{out}}\bm z_t+\bm c)_u\bigr)}.
  \label{eq:nlp-llm-softmax}
\end{equation}
$\theta$包括嵌入、各层参数与输出层，$\bm c\in\mathbb R^V$是输出偏置。输出偏置和参数绑定都是可变的设计：LLaMA-3省去偏置并使用独立的$\bm W_{\mathrm{out}}$，GPT-2则令$\bm W_{\mathrm{out}}=\bm E$。这些选择改变参数化，不改变“有限个得分经softmax归一化为词表分布”这一关系。

式\eqref{eq:nlp-llm-softmax}在数值上需要小心。得分可能达到几十，直接取指数会在低精度格式中溢出。由于分子分母同乘$\exp(-m)$不改变比值，实现中先减去该位置得分的最大值$m$再取指数，并在较高精度下求和与取对数；注意力中逐行的softmax也这样处理。掩码中的$-\infty$在实现中常用该精度下的极小值代替，效果相同：经过指数后权重为零。输出层的计算量也不小，每个位置要算$V$个得分，LLaMA-3-8B中这一步约占每词元参数项计算量的7\%，训练时还要为$B\times T\times V$个得分保存中间结果。

由概率的链式法则，各位置的条件分布连乘得到序列概率。为让完整文本有明确的起止，在序列前加入开始标记$x_0=\mathrm{BOS}$，并把结束标记$x_{T+1}=\mathrm{EOS}$也作为被预测的事件：
\begin{equation}
  p_\theta(x_{1:T})=\prod_{t=1}^{T+1}p_\theta(x_t\mid x_{0:t-1}),
  \qquad x_0=\mathrm{BOS},\ x_{T+1}=\mathrm{EOS}.
  \label{eq:nlp-llm-sequence-probability}
\end{equation}
若不计入EOS，乘积只是“以$x_{1:T}$开头”这一前缀事件的概率，较短的前缀总是概率更大，不同长度的完整文本也就无法比较。计入EOS后，各种长度的完整文本才分享同一个概率总量。

训练与生成以不同的方式使用式\eqref{eq:nlp-llm-sequence-probability}。训练时整条序列已知，输入$x_{0:T}$，目标$x_{1:T+1}$，因果掩码保证位置$t$的输出只依赖$x_{0:t}$，所以一次前向就同时得到全部$T+1$个条件分布，把各位置的负对数概率相加即得损失。按位置平均写为
\begin{equation}
  \mathcal L(\theta)=-\frac{1}{T+1}\sum_{t=1}^{T+1}\log p_\theta(x_t\mid x_{0:t-1}),
  \label{eq:nlp-llm-token-loss}
\end{equation}
它就是第\ref{chap:nlp-classical-models}章的交叉熵，$\exp(\mathcal L)$即该序列上的困惑度。由于分母按词元计数，同一段文本在不同词元化方案下的$\mathcal L$不能直接比较，这一点在第\ref{sec:nlp-llm-tokenization}小节讨论词元化时已经说明。这种在每个位置都以真实前缀作条件的做法称为\term{教师强制}（teacher forcing），它让序列长度上的计算完全并行。生成时则没有真实前缀：必须先确定$x_t$，才能计算$x_{t+1}$的分布，各步只能依次进行。

单个词元的概率与整句的概率也不是一回事。用一个教学例子说明：第一步$p(A)=0.6$、$p(B)=0.4$；$A$之后有两个续写各占0.5，$B$之后有一个续写占0.9。每步选最可能的词元会先选$A$，得到的两词序列概率为$0.6\times0.5=0.30$；而从$B$出发的序列概率为$0.4\times0.9=0.36$。单步最优的选择不保证得到概率最高的序列。怎样在生成时近似地寻找高概率序列、何时改用采样、怎样复用已经算出的键和值，是第\ref{chap:nlp-model-inference}章的内容。到这里，从词元编号到序列概率的计算已经完整写出；第\ref{sec:nlp-llm-long-context}节从这条基线出发，处理注意力在长上下文下的计算量与状态量。
